\subsection*{DMA 语句族}

分块搬运使用统一的基本形式：

\begin{lstlisting}[language=croqtile,caption={DMA 的一般形式}]
future = engine.operation.modifiers source => destination;
\end{lstlisting}

各部分含义如下：

\begin{table}[htbp]\centering\small
\begin{tabularx}{\linewidth}{@{}lX@{}}
\toprule
组成部分 & 含义 \\
\midrule
\co@future@ & 可选的传输结果句柄。用 \co@future.data@ 访问已搬运分块，用 \co@future.span@ 查询其形状。 \\
\co@operation@ & \co@copy@、\co@transp<...>@ 或 \co@pad<...>@。 \\
\co@modifiers@ & 可选标志，如 \co@.async@、\co@.zfill@ 或 \co@.swiz<N>@。 \\
\co@source => destination@ & 源区域与目的存储位置或区域。 \\
\bottomrule
\end{tabularx}
\end{table}

最简单的拷贝示例如下：

\begin{lstlisting}[language=croqtile,caption={一个具体的拷贝}]
f = dma.copy input.chunkat(blk) => shared;
g = dma.transp<1, 0> matrix => local;
h = dma.pad<{2, 1}, {3, 2}, {0, 0}, 0> tile => shared;
\end{lstlisting}

搬运对象是带形状的数据区域，可以一次描述整个分块的传输。\co@=>@ 左侧写源区域，右侧写目的区域或存储位置；目的只写 \co@shared@ 时，编译器根据源形状准备缓冲区，后续代码可以通过期值引用它。
